\section{Introduction}
Machine learning models are increasingly distributed through shared registries, fine-tuned by downstream adopters, compressed for edge deployment, and integrated into regulated workflows. In this setting, establishing \emph{model provenance}---connecting a suspicious or modified checkpoint back to a specific developer and a committed training dataset---presents two coupled security challenges. First, an external signature file shipped alongside model weights can be trivially deleted or substituted; exact file hashing breaks as soon as benign lifecycle operations such as quantization, pruning, or fine-tuning alter parameter values. Second, high-value models and their provenance records can remain in service for years, exposing classical digital signatures (RSA, ECDSA, Ed25519) to a ``harvest-now, forge-later'' (HNFL) threat once cryptographically relevant quantum computers execute Shor's algorithm~\cite{shor1994algorithms,campbell2026mbom}.

To resist quantum forgery, NIST standardized the Module-Lattice-Based Digital Signature Algorithm (ML-DSA) in FIPS~204~\cite{nist204}. Yet adopting post-quantum signatures creates an immediate engineering constraint for weight watermarking: an ML-DSA-65 signature occupies 3,309 bytes (26,472 bits), far too large and brittle to embed bit-by-bit inside a compact neural network without destroying accuracy or collapsing upon a single bit flip. Binding the signed record to the weights instead requires decoupling the external post-quantum signature from a compact, error-corrected cryptographic fingerprint embedded inside the model parameters.

Even when a compact fingerprint is embedded via keyed linear projections~\cite{uchida2017embedding}, two critical questions govern practical security. First, white-box watermarks are vulnerable to \emph{neural structural obfuscation}~\cite{yan2023rethinking}: permuting or positively rescaling the channels feeding a carrier layer leaves network predictions mathematically unchanged while scrambling linear weight projections. Second, enforcing high-margin binary cross-entropy (BCE) constraints across hundreds of encoded bits degrades clean classification accuracy by inflating a shared component across carrier filters. Capping the watermark objective with a hinge margin reduces that accuracy penalty, but raises a central experimental question: \emph{does recovering clean classification accuracy via a margin-capped objective preserve resistance to adversarial removal?}

We investigate this question end to end on a 160,234-parameter convolutional classifier (\textsc{SmallCNN}) trained on CIFAR-10~\cite{cifar10}. Our contributions are fourfold:
\begin{itemize}
  \item \textbf{Post-quantum provenance pipeline:} We combine a domain-separated SHA-256 Merkle commitment over all 50,000 CIFAR-10 training images, a canonical JSON release manifest signed with ML-DSA-65, and a soft-decoded Reed--Solomon (RS32/RS16) keyed projection watermark carrying the full 256-bit manifest digest~\cite{projectResults}.
  \item \textbf{Provably bounded structural-obfuscation recovery:} We introduce a white-box verifier that reverses function-preserving input-channel permutations (via Hungarian assignment) and per-channel positive scaling (via linear programming) strictly when Cover's function-counting theorem~\cite{cover1965geometrical} guarantees a false-fit bound below $2^{-64}$, alongside a binomial residual-evidence tier (\textsc{Damaged}) for overwrite forensics.
  \item \textbf{Mechanistic accuracy--robustness analysis:} Across five paired training seeds, we show that a margin-capped hinge loss ($m=1.5$) reduces shared-filter norm inflation by more than half and recovers $+1.73$ percentage points of clean accuracy over BCE/RS32 (84.10\% vs.\ 82.37\%).
  \item \textbf{Utility-conditioned robustness benchmark:} Across 252 checkpoint evaluations per design covering 47 modifications (noise, pruning, quantization, fine-tuning, channel permutation/scaling, and watermark overwrite) and 778 negative controls, we show that the hinge objective fails all three 512-bit overwrite trials within a 5-point utility budget, whereas BCE/RS32 and BCE/RS16 incur zero successful removals and zero false positives.
\end{itemize}
